\documentclass[12pt,a4paper]{article}
\usepackage[romanian]{babel}
\usepackage[utf8]{inputenc}
\usepackage[T1]{fontenc}
\usepackage{geometry}
\usepackage{graphicx}
\usepackage{hyperref}
\usepackage{float}
\usepackage{amsmath}
\usepackage{listings}

\geometry{margin=2.5cm}

\lstset{
    basicstyle=\ttfamily\small,
    breaklines=true,
    columns=fullflexible
}

\begin{document}
\begin{titlepage}
\begin{center}
\vspace*{1cm}
{\Large Technical University of Moldova}\\[0.5cm]
{\Large Faculty of Computers, Informatics and Microelectronics}\\[2cm]
{\Huge\bfseries Laboratory Work No. 4.1}\\[1cm]
{\Large\bfseries Embedded Systems:\\
Binary Signal Acquisition with Threshold Alerting\\
and Periodic STDIO Reporting}\\[3.5cm]
\begin{flushright}
\begin{tabular}{ll}
\textbf{Elaborated:} & \\
std. gr. FAF-231 & Burduja Adrian \\[1cm]
\textbf{Verified:} & \\
asist. univ. & Alexei Martîniuc \\ & \\
\end{tabular}
\end{flushright}
\vfill
{\Large Chișinău -- 2026}
\end{center}
\end{titlepage}
\newpage

%-------------------------------------------------
\section{Task Analysis and Technological Context}

\subsection{Domain Analysis and Technological Context}
Binary signal conditioning is a fundamental pattern in embedded monitoring systems.
Rather than treating a sensor reading as a continuous value, the system maps it to a
binary alert state: normal or alarm. This pattern appears in thermostats, industrial
safety monitors, and environmental control systems.

A naive implementation that triggers an alert the moment a threshold is crossed suffers
from oscillation — when the measured value hovers near the threshold, the alert state
flips rapidly. Two classical techniques address this: hysteresis (using separate activation
and deactivation levels) and persistence debouncing (requiring the condition to hold for
several consecutive samples before committing a state change).

This work implements binary threshold alerting for a DHT22 temperature/humidity sensor
and an HC-SR04 ultrasonic distance sensor on an Arduino Mega, with structured STDIO reporting.

\subsection{External Resources and Bibliographic References}
\begin{itemize}
    \item documentation of the Arduino Mega (ATmega2560) microcontroller;
    \item DHT22 sensor datasheet;
    \item HC-SR04 ultrasonic sensor datasheet;
    \item AVR-libc documentation (\texttt{dtostrf}, \texttt{fdev\_setup\_stream});
    \item educational materials provided within the course;
    \item Claude (Anthropic) for assistance with code generation and debugging.
\end{itemize}

\subsection{Analysis of Components and Used Resources}
\begin{itemize}
    \item Arduino Mega (ATmega2560) microcontroller;
    \item DHT22 sensor on pin A0 — temperature and humidity;
    \item HC-SR04 ultrasonic sensor on TRIG=7, ECHO=6 — distance;
    \item STDIO interface via serial port at 9600 baud;
    \item software resources: hysteresis logic, persistence counter, millis-based scheduling.
\end{itemize}

%-------------------------------------------------
\section{Definition of Technical Requirements and Test Scenarios}

\subsection{Application Objectives}
\begin{itemize}
    \item periodically acquire raw data from both sensors;
    \item apply hysteresis-based threshold detection to temperature and distance;
    \item debounce state transitions using a persistence counter;
    \item report sensor values and alert states via STDIO at a configurable interval.
\end{itemize}

\subsection{Functional and Non-Functional Requirements}

\textbf{Functional requirements:}
\begin{itemize}
    \item acquisition of temperature, humidity, and distance at \texttt{ACQUIRE\_INTERVAL};
    \item temperature alert above $26$\,°C, cleared below $24$\,°C ($25$\,°C $\pm1$\,°C hysteresis);
    \item distance alert below $18$\,cm, cleared above $22$\,cm ($20$\,cm $\pm2$\,cm hysteresis);
    \item state change requires \texttt{PERSIST\_REQUIRED = 3} consecutive confirming samples;
    \item structured STDIO report at \texttt{REPORT\_INTERVAL}.
\end{itemize}

\textbf{Non-functional requirements:}
\begin{itemize}
    \item modular code in independent \texttt{.cpp}/\texttt{.h} files;
    \item no dynamic memory allocation;
    \item \texttt{printf} via AVR-libc \texttt{fdev\_setup\_stream};
    \item float formatting via \texttt{dtostrf}.
\end{itemize}

\subsection{Test Scenarios and Validation Criteria}
\begin{itemize}
    \item temperature held above $26$\,°C for 3 samples $\Rightarrow$ ALERT;
    \item temperature held below $24$\,°C for 3 samples $\Rightarrow$ NORMAL;
    \item temperature oscillating in $[24, 26]$\,°C $\Rightarrow$ state unchanged;
    \item single spike above threshold $\Rightarrow$ no state change (persistence blocks it);
    \item distance held below $18$\,cm for 3 samples $\Rightarrow$ TOO CLOSE alert.
\end{itemize}

%-------------------------------------------------
\section{Architectural Design and Behavioral Modeling}

\subsection{System Components}
\begin{itemize}
    \item \textbf{Acquisition} — \texttt{dht\_read()}, \texttt{hcsr\_read()} read raw sensor values;
    \item \textbf{Threshold logic} — \texttt{process\_threshold()}, \texttt{process\_distance()} implement hysteresis and persistence;
    \item \textbf{Reporting} — \texttt{report\_temperature()}, \texttt{report\_distance()} format and print via STDIO;
    \item \textbf{STDIO layer} — \texttt{stdio\_init()} binds \texttt{printf} to the hardware UART.
\end{itemize}

\subsection{HW--SW Layered Architecture}
\begin{itemize}
    \item hardware layer (DHT22, HC-SR04, UART);
    \item HAL layer (GPIO, \texttt{pulseIn}, \texttt{delayMicroseconds});
    \item application layer (threshold logic, persistence, reporting).
\end{itemize}

\begin{figure}[H]
    \centering
    \includegraphics[width=0.8\textwidth]{img/structure.png}
    \caption{Structural diagram}
\end{figure}

\subsection{Hysteresis and Persistence Logic}

Two separate crossing levels prevent alert oscillation near the threshold:

\[
T_{high} = T_{threshold} + T_{hysteresis}, \qquad T_{low} = T_{threshold} - T_{hysteresis}
\]

\begin{itemize}
    \item When \textbf{NORMAL}: alert activates only if value $> T_{high}$;
    \item When \textbf{ALERT}: alert clears only if value $< T_{low}$;
    \item Values in the dead zone $[T_{low},\, T_{high}]$ leave state unchanged.
\end{itemize}

A persistence counter \texttt{g\_persist} increments on each sample that would
cause a transition. The transition is committed only when the counter reaches
\texttt{PERSIST\_REQUIRED}. Any sample that does not continue the trend resets
the counter to zero.

\subsection{Scheduling}
No RTOS is used. Two \texttt{millis()}-based timers control execution:
\begin{itemize}
    \item \texttt{ACQUIRE\_INTERVAL = 2000\,ms} — sensor read and threshold processing;
    \item \texttt{REPORT\_INTERVAL = 500\,ms} — STDIO report.
\end{itemize}
The DHT22 requires a minimum of 2 seconds between reads, which sets the acquisition interval.

\begin{figure}[H]
    \centering
    \includegraphics[width=0.8\textwidth]{img/sequence.png}
    \caption{Sequence diagram}
\end{figure}
%-------------------------------------------------
\section{Electrical Schematic Design and Simulation}

\subsection{Connections and Component Sizing}
\begin{itemize}
    \item DHT22: data pin on A0, $10\,k\Omega$ pull-up resistor to $+5\,V$;
    \item HC-SR04: TRIG on pin 7 (output), ECHO on pin 6 (input), powered at $5\,V$;
    \item serial output via USB at 9600 baud.
\end{itemize}

\subsection{Simulation}
The application was simulated in SimulIDE with virtual sensor stimulus values
to verify hysteresis boundaries and persistence debouncing behavior.

\begin{figure}[H]
    \centering
    \includegraphics[width=0.8\textwidth]{img/sim.PNG}
    \caption{Electrical schematic in simulator}
\end{figure}

\subsection{Execution on Real Hardware}
Sample serial output:
\begin{lstlisting}
=== DHT22 Temperature Monitor (Mega) ===
------------------------------
Uptime (ms) : 2001
Temperature :  24.5 C
Humidity    :  58.0 %
Threshold   :  25.0 C  +/- 1.0 C
Persist     : 0/3
State       : NORMAL
Distance    :  35.2 cm
Dist Thresh :  20.0 cm
Dist State  : NORMAL
------------------------------
Uptime (ms) : 4002
Temperature :  26.3 C
Humidity    :  57.5 %
Threshold   :  25.0 C  +/- 1.0 C
Persist     : 1/3
State       : NORMAL
------------------------------
Uptime (ms) : 6003
Temperature :  26.8 C
Humidity    :  57.2 %
Threshold   :  25.0 C  +/- 1.0 C
Persist     : 2/3
State       : NORMAL
------------------------------
Uptime (ms) : 8004
Temperature :  27.1 C
Humidity    :  57.0 %
Threshold   :  25.0 C  +/- 1.0 C
Persist     : 0/3
State       : !!! ALERT !!!
\end{lstlisting}

%-------------------------------------------------
\section{Modular Software Implementation}

\subsection{Project Structure}
\begin{lstlisting}
lab4-1.ino
main.cpp
dht22.cpp
dht22.h
hc-sr.cpp
hc-sr.h
stdio_layer.cpp
stdio_layer.h
\end{lstlisting}

\subsection{HW--SW Interface}
Sensor access is encapsulated behind \texttt{dht\_read()} and \texttt{hcsr\_read()}.
Threshold logic is encapsulated in \texttt{process\_threshold()} and \texttt{process\_distance()},
operating on module-level global state exposed via \texttt{extern} declarations in headers.

\subsection{Key Implementation Details}
\begin{itemize}
    \item \texttt{LOW} and \texttt{HIGH} are macros in \texttt{Arduino.h} (expanding to 0 and 1);
          threshold variables were named \texttt{THRESH\_LOW} / \texttt{THRESH\_HIGH} to avoid
          silent macro substitution corrupting float comparisons;
    \item \texttt{printf} \texttt{\%f} is unsupported on AVR-libc without special linker flags;
          \texttt{dtostrf()} is used for all float formatting;
    \item \texttt{pulseIn} with a 30\,ms timeout prevents blocking indefinitely on out-of-range
          HC-SR04 readings; a return value of $-1$ signals an invalid measurement.
\end{itemize}

%-------------------------------------------------
\section{Simulation, Testing, and Validation}

\subsection{Testing}
Tests were performed by providing controlled sensor stimulus in SimulIDE and
observing serial output to verify hysteresis boundaries, persistence counting,
and correct alert state transitions.

\subsection{Results and Observations}
\begin{itemize}
    \item hysteresis correctly prevented alert oscillation when temperature hovered near $25$\,°C;
    \item persistence counter blocked false alerts from single transient spikes;
    \item both sensors reported independently without interfering with each other's state;
    \item STDIO output was correctly formatted for all measured values.
\end{itemize}

%-------------------------------------------------
\section{Conclusions}

The developed application demonstrates effective binary signal conditioning on a
resource-constrained AVR microcontroller. Hysteresis and persistence debouncing together
eliminate the two main sources of false alerting near a threshold boundary. The modular
architecture separates acquisition, threshold logic, and reporting into independent files,
making the system straightforward to extend with additional sensors or thresholds.

\newpage
\section{Annex}

\subsection{main.cpp}
\begin{lstlisting}
#include <Arduino.h>
#include "stdio_layer.h"
#include "dht22.h"
#include "hc-sr.h"

int main() {
    init();

    stdio_init();

    dht.begin();
    hcsr_init();

    printf("=== DHT22 Temperature Monitor (Mega) ===\r\n");

    unsigned long lastAcquire = 0;
    unsigned long lastReport  = 0;

    for (;;) {
        unsigned long now = millis();

        if (now - lastAcquire >= ACQUIRE_INTERVAL) {
            lastAcquire = now;
            RawData data = dht_read();
            float dist = hcsr_read();
            if (data.valid) { process_threshold(); }
            if (dist >= 0) { process_distance(); }
        }

        if (now - lastReport >= REPORT_INTERVAL) {
            lastReport = now;
            report_temperature(now);
            report_distance();
        }


    }

    return 0;
}
\end{lstlisting}

\subsection{dht22.cpp}
\begin{lstlisting}
#include "dht22.h"
#include <Arduino.h>

DHT dht(DHT_PIN, DHT_TYPE);


RawData    g_data       = { 0, 0, false };
bool       g_is_alert   = false;
int        g_persist    = 0;


RawData dht_read() {
    g_data.temperature = dht.readTemperature();
    g_data.humidity    = dht.readHumidity();
    g_data.valid       = !isnan(g_data.temperature) && !isnan(g_data.humidity);
    return g_data;
}

void process_threshold() {
    float temp = g_data.temperature;
    const float high_thresh = TEMP_THRESHOLD + TEMP_HYSTERESIS;
    const float low_thresh  = TEMP_THRESHOLD - TEMP_HYSTERESIS;

    if (g_is_alert == false) {
        if (temp > high_thresh) {
            if (++g_persist >= PERSIST_REQUIRED) { g_is_alert = true;  g_persist = 0; }
        } else {
            g_persist = 0;
        }
    } else {
        if (temp < low_thresh) {
            if (++g_persist >= PERSIST_REQUIRED) { g_is_alert = false; g_persist = 0; }
        } else {
            g_persist = 0;
        }
    }
}

void report_temperature(unsigned long now) {
    char buf[8];

    printf("------------------------------\r\n");
    printf("Uptime (ms) : %lu\r\n", now);

    if (!g_data.valid) {
        printf("Sensor      : READ ERROR\r\n");
        return;
    }

    printf("Temperature : %s C\r\n",         dtostrf(g_data.temperature, 5, 1, buf));
    printf("Humidity    : %s %%\r\n",         dtostrf(g_data.humidity,    5, 1, buf));
    printf("Threshold   : %s C  +/-",         dtostrf(TEMP_THRESHOLD,     5, 1, buf));
    printf("%s C\r\n",                         dtostrf(TEMP_HYSTERESIS,    4, 1, buf));
    printf("Persist     : %d/%d\r\n",          g_persist, PERSIST_REQUIRED);
    printf("State       : %s\r\n",             g_is_alert == true ? "!!! ALERT !!!" : "NORMAL");
}

\end{lstlisting}

\subsection{dht22.h}
\begin{lstlisting}
#pragma once
#include <DHT.h>

#define DHT_PIN            A0
#define DHT_TYPE           DHT22

#define ACQUIRE_INTERVAL   2000UL
#define REPORT_INTERVAL    500UL

#define TEMP_THRESHOLD     25.0f
#define TEMP_HYSTERESIS     1.0f
#define PERSIST_REQUIRED    3

struct RawData {
    float temperature;
    float humidity;
    bool  valid;
};

extern DHT dht; 

extern RawData    g_data       ;
extern bool       g_is_alert   ;
extern int        g_persist    ;


RawData dht_read();
void process_threshold();
void report_temperature(unsigned long now);
\end{lstlisting}

\subsection{dht22.h}
\begin{lstlisting}
#include "hc-sr.h"
#include <Arduino.h>

float g_distance    = 0.0f;
 bool g_dist_alert  = false;
  int g_dist_persist= 0;

// for string formatting
char buf[8];

float hcsr_read(){
    digitalWrite(TRIG_PIN, LOW);
    delayMicroseconds(2);
    digitalWrite(TRIG_PIN, HIGH);
    delayMicroseconds(10);
    digitalWrite(TRIG_PIN, LOW);

    unsigned long duration = pulseIn(ECHO_PIN, HIGH, 30000UL); // 30ms timeout ~= 5m max
    g_distance = duration / 58.0f;
    if (duration == 0) g_distance -= 1.0f;                  // timeout = out of range
    return g_distance;
}

void process_distance(){
    float dist = g_distance;
    if (dist < 0) { g_dist_persist = 0; return; }     // ignore invalid reads

    const float THRESH_NEAR = DIST_THRESHOLD - DIST_HYSTERESIS;
    const float THRESH_FAR  = DIST_THRESHOLD + DIST_HYSTERESIS;

    if (dist < THRESH_NEAR && !g_dist_alert) {
        if (++g_dist_persist >= DIST_PERSIST_REQUIRED) {
            g_dist_alert  = true;
            g_dist_persist = 0;
        }
    } else if (dist > THRESH_FAR && g_dist_alert) {
        if (++g_dist_persist >= DIST_PERSIST_REQUIRED) {
            g_dist_alert  = false;
            g_dist_persist = 0;
        }
    } else {
        g_dist_persist = 0;
    } 
}

void hcsr_init(){
    pinMode(TRIG_PIN, OUTPUT);
    pinMode(ECHO_PIN, INPUT);
}

void report_distance(){

    printf("Distance    : %s cm\r\n",  dtostrf(g_distance,     5, 1, buf));
    printf("Dist Thresh : %s cm\r\n",  dtostrf(DIST_THRESHOLD, 5, 1, buf));
    printf("Dist State  : %s\r\n",     g_dist_alert ? "!!! TOO CLOSE !!!" : "NORMAL");
}
\end{lstlisting}

\subsection{hc-sr.h}
\begin{lstlisting}
#pragma once

#define TRIG_PIN              7
#define ECHO_PIN              6

#define DIST_THRESHOLD       20.0f   // cm
#define DIST_HYSTERESIS       2.0f   // +-cm  -> alert below 18, clear above 22
#define DIST_PERSIST_REQUIRED 3

extern float g_distance    ;
extern bool  g_dist_alert  ;
extern  int  g_dist_persist;

float hcsr_read();
void process_distance();
void report_distance();

void hcsr_init();

\end{lstlisting}

\subsection{stdio\_layer.cpp}
\begin{lstlisting}
#include "stdio_layer.h"
#include "HardwareSerial.h"
#include <Arduino.h>

static int serial_putchar(char c, FILE *) {
    Serial.write(c);
    return 0;
}

static FILE g_stdout;

void stdio_init() {
    fdev_setup_stream(&g_stdout, serial_putchar, NULL, _FDEV_SETUP_WRITE);
    stdout = &g_stdout;
    Serial.begin(9601);
}
\end{lstlisting}

\subsection{stdio\_layer.h}
\begin{lstlisting}
#pragma once
#include <stdio.h>

void stdio_init();
\end{lstlisting}
\end{document}